Denoising diffusion models~\cite{ho2020denoising,song2020score} and flow matching~\cite{lipman2022flow,liu2022flow,albergo2022building} both turn Gaussian noise into data by integrating a learned vector field, yet they are trained on different paths. Diffusion models regress the noise added along a variance-preserving schedule, and the resulting probability-flow trajectories are typically curved, so deterministic samplers such as DDIM~\cite{song2020denoising} need tens of steps. Flow matching regresses the velocity of a straight line between a noise sample and a data sample. The common argument is that straighter marginal paths can be integrated accurately with few Euler steps, and that rectified flow~\cite{liu2022flow} can straighten them further by \emph{reflow}: retraining on (noise, sample) couplings produced by the model itself.

The published evidence for this argument is mostly on image benchmarks. Lipman et al.~\cite{lipman2022flow} compare flow matching with diffusion-path and score-matching baselines on CIFAR-10 and ImageNet using one shared architecture, and report likelihood, FID and the number of function evaluations of an adaptive solver; Liu et al.~\cite{liu2022flow} report FID and recall of one-step rectified flows on CIFAR-10. Both papers, and minibatch-OT flow matching~\cite{tong2023improving}, also show 2D toy examples. We do not question that evidence. We run a deliberately small, controlled replication on three 2D densities (eight Gaussians, two moons, checkerboard), a setting in which fresh samples of the target are available, so that distributional distances to held-out data can be computed directly, every step count from 1 to 100 can be swept, a DDIM sampler can serve as the diffusion baseline, and seed-to-seed variance can be reported. Every system uses the same MLP, optimiser, number of updates and training data stream, and is scored with the same held-out reference, a sliced Wasserstein distance~\cite{bonneel2015sliced,kolouri2019generalized} and an MMD, at 1 to 100 sampling steps. We ask: (H1) does flow matching give better samples than DDIM at four steps; (H2) does that advantage vanish with more steps; (H3) does one reflow round help one-step sampling without helping at 100 steps; (H4) are flow trajectories actually straighter, as measured by chord-to-arc length.

Contributions: a single-file CPU reimplementation of DDPM/DDIM, flow matching and reflow on a shared network; five-seed step-count curves with a real-data floor; an equal-evaluation comparison of Heun and Euler sampling; ablations on the diffusion schedule, the reflow teacher and a second round; and a hypothesis-by-hypothesis report including failures. The study is small (2D, one MLP, 3000 updates) and its conclusions are limited accordingly (Sec.~\ref{sec:limitations}).
